\subsection{INTEGRATION FORMULA OF H. WEYL}

Let e be the identity element of G and \(\bar{e}\) its class in G/T. Identify the tangent space of G at e with g, the tangent space of T at e with t and the tangent space of G/T at \(\bar{e}\) with g/t. Denote by \((u,v)\mapsto u\cap v\) the R-bilinear map

\[
\operatorname{Alt} ^ {n - r} (\mathfrak {g} / \mathfrak {t}) \times \operatorname{Alt} ^ {r} (\mathfrak {t}) \rightarrow \operatorname{Alt} ^ {n} (\mathfrak {g})
\]

defined in number 1.

Recall (Chap. III, §3, no. 13, Prop. 50) that the map \(\omega \mapsto \omega(e)\) is an isomorphism from the vector space of left-invariant differential forms of degree \(n\) (resp. \(r\) ) on G (resp. T) to the space \(\mathrm{Alt}^n(\mathfrak{g})\) (resp. \(\mathrm{Alt}^r(t)\) ). Further, observe that, since every connected compact subgroup of \(\mathbf{R}^*\) reduces to the identity element, \(\det \mathrm{Ad}g = 1\) for all \(g \in \mathrm{G}\) , so that the left-invariant differential forms of degree \(n\) on G are also right-invariant and invariant under inner automorphisms (Chap. III, §3, no. 16, Cor. of Prop. 54): we shall speak simply of G-invariant differential forms from now on.

Similarly, it follows from Chap. III, §3, no. 16, Prop. 56 and the preceding that the map \(\omega \mapsto \omega (\bar{e})\) is an isomorphism from the space of G-invariant differential forms of degree \(n - r\) on \(\mathrm{G / T}\) to the space \(\mathrm{Alt}^{n - r}(\mathfrak{g} / \mathfrak{t})\) .

If \(\omega_{G/T}\) is a G-invariant differential form of degree n - r on G/T, and \(\omega_{T}\) an invariant differential form of degree r on T, denote by \(\omega_{G/T} \cap \omega_{T}\) the unique invariant differential form of degree n on G such that

\[
(\omega_ {\mathrm{G/T}} \cap \omega_ {\mathrm{T}}) (e) = \omega_ {\mathrm{G/T}} (\bar {e}) \cap \omega_ {\mathrm{T}} (e).
\]

Recall finally that \(f : (\mathrm{G}/\mathrm{T}) \times \mathrm{T} \to \mathrm{G}\) denotes the morphism of manifolds induced by the map \((g, t) \mapsto gtg^{-1}\) from \(G \times T\) to G by passage to the quotient (§5, no. 4). If \(\alpha\) and \(\beta\) are differential forms on G/T and T, respectively, denote simply by \(\alpha \wedge \beta\) the form \(pr_{1}^{*}\alpha \wedge pr_{2}^{*}\beta\) on \((\mathrm{G}/\mathrm{T}) \times \mathrm{T}\) .

For \(t \in \mathrm{T}\) , denote by \(\operatorname{Ad}_{\mathfrak{g} / \mathfrak{t}}(t)\) the endomorphism of \(\mathfrak{g} / \mathfrak{t}\) induced by \(\operatorname{Ad} t\) by passage to the quotient. Put

\[
\delta_ {\mathrm{G}} (t) = \det (\mathrm{Ad} _ {\mathfrak {g} / \mathfrak {t}} (t) - 1) = \prod_ {\alpha \in \mathrm{R} (\mathrm{G}, \mathrm{T})} (t ^ {\alpha} - 1).\tag{3}
\]

Let \(x \in \mathfrak{t}\) and \(\alpha \in \mathrm{R}(\mathrm{G},\mathrm{T})\) ; denote by \(\hat{\alpha}\) the element \((2\pi i)^{-1}\delta (\alpha)\) of \(\mathfrak{t}^*\) , so that

\[
((\exp x) ^ {\alpha} - 1) ((\exp x) ^ {- \alpha} - 1) = (e ^ {2 \pi i \hat {\alpha} (x)} - 1) (e ^ {- 2 \pi i \hat {\alpha} (x)} - 1) = 4 \sin^ {2} \pi \hat {\alpha} (x).
\]

If \(\mathrm{R}_{+}(\mathrm{G},\mathrm{T})\) denotes the set of positive roots of \(\mathrm{R}(\mathrm{G},\mathrm{T})\) relative to a basis B, we have

\[
\delta_ {\mathrm{G}} (\exp x) = \prod_ {\alpha \in \mathrm{R} _ {+} (\mathrm{G}, \mathrm{T})} 4 \sin^ {2} \pi \hat {\alpha} (x),
\]

so, in particular, \(\delta_{\mathrm{G}}(t)>0\) for all \(t\in T_{r}\) . We remark also that \(\delta_{\mathrm{G}}(t)=\delta_{\mathrm{G}}(t^{-1})\) for \(t\in T\) .

PROPOSITION 1. Let \(\omega_{\mathrm{G}}, \omega_{\mathrm{G/T}}\) and \(\omega_{\mathrm{T}}\) be invariant differential forms on G, G/T and T, respectively, of respective degrees \(n, n - r\) and \(r\) . If \(\omega_{\mathrm{G}} = \omega_{\mathrm{G/T}} \cap \omega_{\mathrm{T}}\) , then

\[
f ^ {*} (\omega_ {\mathrm{G}}) = \omega_ {\mathrm{G/T}} \wedge \delta_ {\mathrm{G}} \omega_ {\mathrm{T}}.
\]

Clearly we can assume that \(\omega_{G/T}\) and \(\omega_{T}\) are non-zero; then the differential form \((u,t)\mapsto\omega_{\mathrm{G/T}}(u)\wedge\omega_{\mathrm{T}}(t)\) on \((\mathrm{G/T})\times\mathrm{T}\) is of degree n and everywhere non-zero; hence there exists a numerical function \(\delta\) on \((\mathrm{G/T})\times\mathrm{T}\) such that

\[
f ^ {*} (\omega_ {\mathrm{G}}) (u, t) = \delta (u, t) \omega_ {\mathrm{G/T}} (u) \wedge \omega_ {\mathrm{T}} (t).
\]

Observe now that, for \(h \in \mathrm{G}\) , \(u \in \mathrm{G} / \mathrm{T}\) , \(t \in \mathrm{T}\) , we have \(f(h.u,t) = (\operatorname{Int} h)f(u,t)\) ; since \(\omega_{\mathrm{G}}\) is invariant under inner automorphisms, it follows immediately that \(\delta(h.u,t) = \delta(u,t)\) , so \(\delta(u,t) = \delta(\bar{e},t)\) .

Denote by \(p:\mathfrak{g}\to \mathfrak{g} / \mathfrak{t}\) the quotient map and by \(\varphi :\mathfrak{g} / \mathfrak{t}\rightarrow \mathfrak{g}\) the map defined by

\[
\varphi (p (X)) = (\operatorname{Ad} t ^ {- 1}) X - X \text { for } X \in \mathfrak {g};
\]

recall (§5, no. 4, Lemma 4) that the tangent map

\[
\mathrm{T} _ {(e, t)} (f): \mathrm{T} _ {e} (\mathrm{G} / \mathrm{T}) \times \mathrm{T} _ {t} (\mathrm{T}) \to \mathrm{T} _ {t} (\mathrm{G})
\]

takes \((z,tH)\) to \(t(\varphi (z) + H)\) for \(z\in \mathfrak{g} / \mathfrak{t},H\in \mathfrak{t}\) .

Let \(z_{1},\ldots,z_{n-r}\) be elements of g/t, \(H_{1},\ldots,H_{r}\) elements of t. Then

\[
\begin{array}{l} f ^ {*} \omega_ {\mathrm{G}} (\bar {e}, t) (z _ {1}, \ldots , z _ {n - r}, t H _ {1}, \ldots , t H _ {r}) \\ \quad = \omega_ {\mathrm{G}} (t) (t \varphi (z _ {1}), \ldots , t \varphi (z _ {n - r}), t H _ {1}, \ldots , t H _ {r}) \text {by the calculation of T} _ {(\bar {e}, t)} (f) \\ \quad = \omega_ {\mathrm{G}} (e) (\varphi (z _ {1}), \ldots , \varphi (z _ {n - r}), H _ {1}, \ldots , H _ {r}) \text {since} \omega_ {\mathrm{G}} \text {is invariant} \\ \quad = \omega_ {\mathrm{G/T}} (\bar {e}) (p \varphi (z _ {1}), \ldots , p \varphi (z _ {n - r})) \cdot \omega_ {\mathrm{T}} (e) (H _ {1}, \ldots , H _ {r}) \text {(no. 1, Lemma 1)} \\ \quad = \det (p \varphi) \omega_ {\mathrm{G/T}} (\bar {e}) (z _ {1}, \ldots , z _ {n - r}). \omega_ {\mathrm{T}} (e) (H _ {1}, \ldots , H _ {r}) \\ \quad = \delta_ {\mathrm{G}} (t) \omega_ {\mathrm{G/T}} (\bar {e}) (z _ {1}, \ldots , z _ {n - r}). \omega_ {\mathrm{T}} (t) (t H _ {1}, \ldots , t H _ {r}) \\ \qquad \qquad \qquad \qquad \qquad \qquad \text {since} \omega_ {\mathrm{T}} \text {is invariant} \\ \quad = \delta_ {\mathrm{G}} (t) (\omega_ {\mathrm{G/T}} \wedge \omega_ {\mathrm{T}}) (\bar {e}, t) (z _ {1}, \ldots , z _ {n - r}, t H _ {1}, \ldots , t H _ {r}), \end{array}
\]

so \(f^{*}\omega_{\mathrm{G}}(\bar{e},t) = \delta_{\mathrm{G}}(t)(\omega_{\mathrm{G / T}}\wedge \omega_{\mathrm{T}})(\bar{e},t)\) ; thus, \(\delta (\bar {e},t) = \delta_{\mathrm{G}}(t)\) , hence the proposition.

Give the manifolds G, T and G/T the orientations defined by the forms \(\omega_{G}, \omega_{T}\) and \(\omega_{G/T}\) , respectively. These forms define invariant measures on G, T and G/T (Chap. III, §3, no. 16, Props. 55 and 56), also denoted by \(\omega_{G}, \omega_{T}\) and \(\omega_{G/T}\) .

Lemma 2. If \(\omega_{\mathrm{G}} = \omega_{\mathrm{G / T}}\cap \omega_{\mathrm{T}}\) , then

\[
\int_ {\mathrm{G}} \omega_ {\mathrm{G}} = \int_ {\mathrm{G/T}} \omega_ {\mathrm{G/T}}. \int_ {\mathrm{T}} \omega_ {\mathrm{T}}.
\]

Denote by \(\pi\) the canonical morphism from G to G/T. Let \(g \in G\) , and let \(t_{1}, \ldots, t_{n-r}\) be elements of \(\mathrm{T}_{\pi(g)}(\mathrm{G}/\mathrm{T})\) . Identify the fibre \(\pi^{-1}(\pi(g)) = g\mathrm{T}\) with T by the translation \(\gamma(g)\) . The relation \(\omega_{G} = \omega_{G/T} \cap \omega_{T}\) now implies the equality (Differentiable and Analytic Manifolds, Results, 11.4.5):

\[
\omega_ {\mathrm{G} \sqcup} (t _ {1}, \dots , t _ {n - r}) = (\omega_ {\mathrm{G} / \mathrm{T}} (t _ {1}, \dots , t _ {n - r})) \omega_ {\mathrm{T}}.
\]

Thus \(\int_{\pi}\omega_{\mathrm{G}} = \left(\int_{\mathrm{T}}\omega_{\mathrm{T}}\right)\omega_{\mathrm{G / T}}\) (Differentiable and Analytic Manifolds, Results, 11.4.6), and

\[
\int_ {\mathrm{G}} \omega_ {\mathrm{G}} = \int_ {\mathrm{G/T}} \int_ {\pi} \omega_ {\mathrm{G}} = \int_ {\mathrm{T}} \omega_ {\mathrm{T}}. \int_ {\mathrm{G/T}} \omega_ {\mathrm{G/T}}
\]

(Differentiable and Analytic Manifolds, Results, 11.4.8).

Lemma 3. The inverse image on \((\mathrm{G}/\mathrm{T}) \times \mathrm{T}_{r}\) of the measure dg on \(G_{r}\) under the local homeomorphism \(f_{r}\) (Integration, Chap. V, §6, no. 6) is the measure \(\mu \otimes \delta_{G} dt\) , where \(\mu\) is the unique G-invariant measure on G/T of total mass 1.

Choose an invariant differential form \(\omega_{T}\) (resp. \(\omega_{G/T}\) ) on T (resp. G/T) of maximal degree, such that the measure defined by \(\omega_{T}\) (resp. \(\omega_{G/T}\) ) is equal to dt (resp. \(\mu\) ). Put \(\omega_{G} = \omega_{G/T} \cap \omega_{T}\) . Lemma 2 implies that the measure defined by \(\omega_{G}\) is equal to dg. Let U be an open subset of (G/T) × T \(_{r}\) such that \(f_{r}\) induces an isomorphism from U to an open subset V of G \(_{r}\) . Let \(\varphi\) be a continuous function with compact support in V; denote also by \(\varphi\) the extension of \(\varphi\) to G \(_{r}\) which vanishes outside V. We have

\[
\begin{array}{r l} & {\int_ {\mathrm{V}} \varphi d g = \int_ {\mathrm{V}} \varphi \omega_ {\mathrm{G}} = \int_ {\mathrm{U}} (\varphi \circ f _ {r}) f _ {r} ^ {*} (\omega_ {\mathrm{G}})} \\ & {\qquad = \int_ {\mathrm{U}} (\varphi \circ f _ {r}) \omega_ {\mathrm{G/T}} \wedge \delta_ {\mathrm{G}} \omega_ {\mathrm{T}} \quad (\text {Prop. 1})} \\ & {\qquad = \int_ {\mathrm{U}} (\varphi \circ f _ {r}) d \mu . \delta_ {\mathrm{G}} d t,} \end{array}
\]

hence the lemma.

THEOREM 1 (H. Weyl). The measure \(dg\) on \(G\) is the image under the map \((g, t) \mapsto gtg^{-1}\) from \(G \times T\) to \(G\) of the measure \(dg \otimes \frac{1}{w(G)}\delta_G dt\) , where

\[
\delta_ {\mathrm{G}} (t) = \det (\operatorname{Ad} _ {\mathfrak {g} / \mathfrak {t}} (t) - 1) = \prod_ {\alpha \in \mathrm{R} (\mathrm{G}, \mathrm{T})} (t ^ {\alpha} - 1).
\]

Equivalently (Integration, Chap. V, §6, no. 3, Prop. 4), \(dg\) is the image under the map \(f: (\mathrm{G} / \mathrm{T}) \times \mathrm{T} \to \mathrm{G}\) of the measure \(\mu \otimes \frac{1}{w(\mathrm{G})}\delta_{\mathrm{G}}dt\) .

We prove the last assertion. It follows from §5, no. 1 and Differentiable and Analytic Manifolds, Results, 10.1.3 c) that G-G \(_{r}\) is negligible in G and T-T \(_{r}\) is negligible in T. Further, the map \(f_{r}\) makes (G/T) × T \(_{r}\) a principal covering of G \(_{r}\) , with group W (§5, no. 4, Prop. 4 b)). The theorem now follows from Lemma 3 and Integration, Chap. V, §6, no. 6, Prop. 11.

COROLLARY 1. (i) Let \(\varphi\) be an integrable function on G with values in a Banach space or in \(\overline{\mathbf{R}}\) . For almost all \(t \in \mathrm{T}\) , the function \(g \mapsto \varphi(gtg^{-1})\) on G is integrable for \(dg\) . The function \(t \mapsto \delta_{\mathrm{G}}(t) \int_{\mathrm{G}} \varphi(gtg^{-1}) dg\) is integrable on T, and we have

\[
\int_ {\mathrm{G}} \varphi (g) d g = \frac {1}{w (\mathrm{G})} \int_ {\mathrm{T}} \left(\int_ {\mathrm{G}} \varphi (g t g ^ {- 1}) d g\right) \delta_ {\mathrm{G}} (t) d t\tag{4}
\]

(``Hermann Weyl's integration formula'').

\begin{enumerate}
\def\labelenumi{(\roman{enumi})}
\setcounter{enumi}{1}
\tightlist
\item
  Let \(\varphi\) be a positive measurable function on \(\mathbf{G}\) . For almost all \(t \in \mathrm{T}\) , the function \(g \mapsto \varphi(gtg^{-1})\) on \(\mathbf{G}\) is measurable. The function \(t \mapsto \int_{\mathbf{G}}^{*} \varphi(gtg^{-1}) dg\) on \(\mathrm{T}\) is measurable, and we have
\end{enumerate}

\[
\int_ {\mathrm{G}} ^ {*} \varphi (g) d g = \frac {1}{w (\mathrm{G})} \int_ {\mathrm{T}} ^ {*} \left(\int_ {\mathrm{G}} ^ {*} \varphi (g t g ^ {- 1}) d g\right) \delta_ {\mathrm{G}} (t) d t.\tag{5}
\]

Since the map f is induced by passage to the quotient from the map \((g, t) \mapsto gtg^{-1}\) from \(G \times T \to G\) , it suffices to apply Integration, Chap. V, §5, 6, 8 and Integration, Chap. VII, §2.

COROLLARY 2. Let \(\varphi\) be a central function on \(\mathbf{G}\) (that is, such that \(\varphi(gh) = \varphi(hg)\) for all \(g\) and \(h\) in \(\mathbf{G}\) ) with values in a Banach space or in \(\overline{\mathbf{R}}\) .

\begin{enumerate}
\def\labelenumi{\alph{enumi})}
\item
  \(\varphi\) is measurable if and only if its restriction to T is measurable.
\item
  \(\varphi\) is integrable if and only if the function \((\varphi |\mathrm{T})\delta_{\mathrm{G}}\) is integrable on \(\mathrm{T}\) , and in that case we have
\end{enumerate}

\[
\int_ {\mathrm{G}} \varphi (g) d g = \frac {1}{w (\mathrm{G})} \int_ {\mathrm{T}} \varphi (t) \delta_ {\mathrm{G}} (t) d t.\tag{6}
\]

Denote by \(p: \mathrm{G} / \mathrm{T} \times \mathrm{T} \to \mathrm{T}\) the second projection. We have \(\varphi \circ f = (\varphi | \mathrm{T}) \circ p\) ; further, the image under \(p\) of the measure \(\mu \otimes \frac{1}{w(\mathrm{G})} \delta_{\mathrm{G}} dt\) is \(\frac{1}{w(\mathrm{G})} \delta_{\mathrm{G}} dt\) . The corollary now follows from Th. 1 above and Th. 1 of Integration, Chap. V, §6, no. 2, applied to the two proper maps \(f\) and \(p\) .

COROLLARY 3. Let H be a connected closed subgroup of G containing T, h its Lie algebra, and dh the Haar measure on H of total mass 1. Let \(\varphi\) be an integrable central function on G, with values in a Banach space or in R. Then the function \(h \mapsto \varphi(h) \det(\mathrm{Ad}_{\mathfrak{g}/\mathfrak{h}}(h) - 1)\) is integrable and central on H and we have

\[
\int_ {\mathrm{G}} \varphi (g) d g = \frac {w (\mathrm{H})}{w (\mathrm{G})} \int_ {\mathrm{H}} \varphi (h) \mathrm{det} (\mathrm{Ad} _ {\mathfrak {g} / \mathfrak {h}} (h) - 1) d h.\tag{7}
\]

Indeed, the function \(h \mapsto \varphi(h) \det(\mathrm{Ad}_{\mathfrak{g}/\mathfrak{h}}(h) - 1)\) is a central function on H whose restriction to T is the function \(t \mapsto \varphi(t) \delta_{\mathrm{G}}(t) \delta_{\mathrm{H}}(t)^{-1}\) . Thus, the corollary follows from Cor. 2 applied to G and to H.

Remarks. 1) If we take \(\varphi = 1\) in Cor. 3, we obtain

\[
\int_ {\mathrm{H}} \det (\mathrm{Ad} _ {\mathfrak {g} / \mathfrak {h}} (h) - 1) d h = w (\mathrm{G}) / w (\mathrm{H})\tag{8}
\]

and in particular

\[
\int_ {\mathrm{T}} \delta_ {\mathrm{G}} (t) d t = w (\mathrm{G}).\tag{9}
\]

\begin{enumerate}
\def\labelenumi{\arabic{enumi})}
\setcounter{enumi}{1}
\tightlist
\item
  Let \(\nu\) be the measure on the quotient T/W defined by
\end{enumerate}

\[
\int_ {\mathrm{T/W}} \psi (\tau) d \nu (\tau) = \frac {1}{w (\mathrm{G})} \int_ {\mathrm{T}} \psi (\pi (t)) \delta_ {\mathrm{G}} (t) d t,
\]

where \(\pi\) denotes the canonical projection of T onto T/W. Cor. 2 means that the homeomorphism \(T / W \to G / \operatorname{Int}(G)\) (§2, no. 5, Cor. 1 of Prop. 5)

transports the measure \(\nu\) to the image of the measure dg under the canonical projection \(\mathrm{G} \rightarrow \mathrm{G}/\mathrm{Int}(\mathrm{G})\) .

\begin{enumerate}
\def\labelenumi{\arabic{enumi})}
\setcounter{enumi}{2}
\tightlist
\item
  Assume that G is simply-connected. Let A be an alcove of t, and dx the Haar measure on t such that \(\int_{A} dx = 1\) . Then the measure \(\nu\) can also be obtained by transporting the measure \(\frac{1}{w(G)} \prod_{\alpha \in R_{+}(G,T)} 4 \sin^{2} \pi \hat{\alpha}(x) dx\) on \(\overline{A}\) by the homeomorphism \(\overline{A} \rightarrow T/W\) (§5, no. 2, Cor. 1 of Prop. 2).
\end{enumerate}

Example. Take G to be the group \(\mathbf{SU}(2,\mathbf{C})\) and T to be the subgroup of diagonal matrices ( \(\S3\) , no. 6); identify t with R by the choice of basis \(\{iH\}\) of t (loc. cit.). Put A = 10, \(\pi\) ; this is an alcove of t. The interval \(\overline{A} = [0, \pi]\) can be identified with the space of conjugacy classes of G, the element \(\theta\) of \(\overline{A}\) corresponding to the conjugacy class of \(\begin{pmatrix} e^{i\theta} & 0 \\ 0 & e^{-i\theta} \end{pmatrix}\) . Let \(d\theta\) be Lebesgue measure on \([0, \pi]\) ; it follows from the preceding that the image on \(\overline{A}\) of the Haar measure on G is the measure \(\frac{2}{\pi}\sin^{2}\theta d\theta\) .

